\documentclass[12pt,a4paper]{article}

% ============================================================
% CREST-2 Trial - Asymptomatic Carotid Stenosis
% 整理：謝慕揚醫師 MD, PhD, FESC
% ============================================================

% Chinese font setup
\usepackage{xeCJK}
\usepackage{fontspec}
\IfFontExistsTF{Noto Serif CJK TC}{
  \setCJKmainfont{Noto Serif CJK TC}
  \setCJKsansfont{Noto Serif CJK TC}
  \setCJKmonofont{Noto Serif CJK TC}
}{\setCJKmainfont{SimSun}
  \setCJKsansfont{SimSun}
  \setCJKmonofont{SimSun}
}

% Essential packages
\usepackage[top=2.5cm,bottom=2.5cm,left=2cm,right=2cm,headheight=25pt]{geometry}
\usepackage{graphicx}
\usepackage{booktabs}
\usepackage{longtable}
\usepackage{array}
\usepackage{xcolor}
\usepackage{hyperref}
\usepackage{tcolorbox}
\usepackage{titlesec}
\usepackage{enumitem}
\usepackage{amsmath}
\usepackage{multirow}
\usepackage{fancyhdr}

% Color definitions
\definecolor{emergency_red}{RGB}{186,24,27}
\definecolor{emergency_blue}{RGB}{0,114,188}
\definecolor{emergency_gray}{RGB}{120,120,120}
\definecolor{highlight_yellow}{RGB}{255,250,205}

% Hyperlink setup
\hypersetup{
    colorlinks=true,
    linkcolor=emergency_blue,
    citecolor=emergency_blue,
    urlcolor=emergency_blue,
    pdfborder={0 0 0}
}

% Title formats
\titleformat{\section}{\Large\bfseries\color{emergency_red}}{\thesection}{1em}{}
\titleformat{\subsection}{\large\bfseries\color{emergency_blue}}{\thesubsection}{1em}{}
\titleformat{\subsubsection}{\normalsize\bfseries\color{emergency_gray}}{\thesubsubsection}{1em}{}

% Date command
\newcommand{\mydate}{\the\year-\the\month-\the\day}

% Custom header
\pagestyle{fancy}
\fancyhf{}
\fancyhead[L]{\textcolor{emergency_red}{NEJM 教學講義}}
\fancyhead[R]{\textcolor{emergency_blue}{CREST-2: 無症狀頸動脈狹窄治療}}
\fancyfoot[C]{\textcolor{emergency_gray}{\thepage}}
\renewcommand{\headrulewidth}{1pt}
\renewcommand{\headrule}{\hbox to\headwidth{\color{emergency_red}\leaders\hrule height \headrulewidth\hfill}}

\begin{document}

% ============================================================
% Title Page
% ============================================================
\begin{center}
{\LARGE\bfseries\textcolor{emergency_red}{The New England Journal of Medicine}}\\[0.5cm]
{\Large\textbf{Medical Management and Revascularization for Asymptomatic Carotid Stenosis}}\\[0.3cm]
{\large 無症狀頸動脈狹窄的內科治療與血管重建術}\\[0.5cm]
{\normalsize N Engl J Med 2026;394:219-31. DOI: \href{https://doi.org/10.1056/NEJMoa2508800}{10.1056/NEJMoa2508800}}\\[0.3cm]
{\normalsize 整理：謝慕揚 MD, PhD, FESC \quad 日期：\mydate}
\end{center}

\vspace{0.5cm}

% ============================================================
% Key Findings Box
% ============================================================
\begin{tcolorbox}[colback=yellow!10!white,colframe=orange!75!black,title=\textbf{核心發現摘要}]
\textbf{重大發現：}
\begin{itemize}[noitemsep]
    \item Carotid artery stenting (CAS) + 強化內科治療 vs 單獨強化內科治療：4年主要終點事件率 \textbf{2.8\% vs 6.0\%}，顯著減少（P=0.02），NNT=31
    \item Carotid endarterectomy (CEA) + 強化內科治療 vs 單獨強化內科治療：4年主要終點事件率 \textbf{3.7\% vs 5.3\%}，\textcolor{red}{未達統計顯著差異}（P=0.24）
\end{itemize}

\textbf{臨床意義：}對於高度（$\geq$70\%）無症狀頸動脈狹窄患者，CAS 加上強化內科治療可顯著降低中風風險；但 CEA 的額外效益未獲證實。此結果挑戰了傳統以 CEA 為首選的治療觀念。
\end{tcolorbox}

\tableofcontents
\newpage

% ============================================================
\section{研究背景}
% ============================================================

\subsection{臨床問題}

\begin{itemize}
    \item 高度頸動脈狹窄（carotid stenosis）的治療方式在國際間差異極大
    \item 部分國家僅對有症狀患者進行 revascularization，而美國 \textbf{75-80\%} 接受 CAS 或 CEA 的患者為無症狀
    \item 1990-2000年代的隨機試驗顯示 CEA 優於當時的內科治療
    \item 然而，現代內科治療（intensive medical management）、CAS 技術及 CEA 技術均有顯著進步
    \item 近期小型試驗（SPACE-2、ECST-2）挑戰了 revascularization 的額外效益
\end{itemize}

\subsection{研究目標}

\textbf{CREST-2 試驗}旨在評估：對於高度無症狀頸動脈狹窄患者，在強化內科治療的基礎上，\textbf{加入 CAS 或 CEA 是否能提供額外的中風預防效益}。

% ============================================================
\section{研究設計}
% ============================================================

\subsection{試驗架構}

\begin{itemize}
    \item \textbf{設計：}兩個平行的、評估者盲性（observer-blinded）隨機對照試驗
    \item \textbf{研究期間：}2014年12月至2025年（stenting trial 追蹤至2025年7月；CEA trial 至2024年9月）
    \item \textbf{中心數：}155個中心，涵蓋澳洲、加拿大、以色列、西班牙、美國
    \item \textbf{追蹤時間：}中位數 3.6 年（stenting trial）；4.0 年（CEA trial）
    \item \textbf{註冊：}ClinicalTrials.gov NCT02089217
\end{itemize}

\begin{tcolorbox}[colback=emergency_blue!5!white,colframe=emergency_blue,title=\textbf{兩個平行試驗設計}]
\textbf{Stenting Trial (n=1245):}
\begin{itemize}[noitemsep]
    \item 強化內科治療組 (Medical therapy alone, n=629)
    \item CAS + 強化內科治療組 (Stenting group, n=616)
\end{itemize}

\textbf{Endarterectomy Trial (n=1240):}
\begin{itemize}[noitemsep]
    \item 強化內科治療組 (Medical therapy alone, n=623)
    \item CEA + 強化內科治療組 (Endarterectomy group, n=617)
\end{itemize}
\end{tcolorbox}

\subsection{收案與排除條件}

\begin{longtable}{p{3cm}p{11cm}}
\toprule
\textbf{類型} & \textbf{標準} \\
\midrule
\endfirsthead

\toprule
\textbf{類型} & \textbf{標準} \\
\midrule
\endhead

\bottomrule
\endfoot

\textbf{納入條件} &
$\bullet$ 年齡 $\geq$35 歲\newline
$\bullet$ 頸動脈狹窄 $\geq$70\%（Doppler 超音波 peak systolic velocity $\geq$230 cm/sec，且符合以下任一條件：end diastolic velocity $\geq$100 cm/sec、ICA/CCA peak systolic velocity ratio $\geq$4.0、或 CTA/MRA/catheter angiography 顯示 $\geq$70\% 狹窄）\newline
$\bullet$ 隨機分配前180天內無同側頸動脈領域的 stroke、TIA 或 amaurosis fugax \\
\midrule
\textbf{排除條件} &
$\bullet$ 既往 disabling stroke\newline
$\bullet$ 不穩定心絞痛（unstable angina）\newline
$\bullet$ 需要抗凝治療的心房顫動（atrial fibrillation） \\
\bottomrule
\end{longtable}

\subsection{強化內科治療方案}

所有患者均接受相同的強化內科治療方案：

\begin{longtable}{p{4cm}p{10cm}}
\toprule
\textbf{治療目標} & \textbf{內容} \\
\midrule
\endfirsthead

\bottomrule
\endfoot

\textbf{血壓控制} & 收縮壓 <130 mmHg（2018年前目標為 <140 mmHg，後依據指引更新調降） \\
\midrule
\textbf{血脂控制} & LDL cholesterol <70 mg/dL（1.80 mmol/L）\newline
2018年後可提供 alirocumab（PCSK9 inhibitor，由 Regeneron 捐贈） \\
\midrule
\textbf{血糖控制} & 監測並管理 glucose 及 HbA1c \\
\midrule
\textbf{生活型態} & 戒菸、體重控制、規律運動\newline
提供電話健康指導（health coaching） \\
\midrule
\textbf{抗血小板治療} & 依各治療組別給予 aspirin $\pm$ clopidogrel \\
\bottomrule
\end{longtable}

\subsection{介入治療方案}

\subsubsection{Carotid Artery Stenting (CAS)}

\begin{itemize}
    \item 經股動脈入路（transfemoral approach），局部麻醉 $\pm$ conscious sedation
    \item \textbf{必須使用 embolic protection device}（99.6\% 使用率）
    \item 術前48小時開始：Aspirin 325 mg QD + Clopidogrel 75 mg BID
    \item 術後30天：Clopidogrel 75 mg QD + Aspirin 75-325 mg QD
    \item 30天後：Aspirin 70-325 mg QD 持續使用
\end{itemize}

\subsubsection{Carotid Endarterectomy (CEA)}

\begin{itemize}
    \item 89\% 使用全身麻醉（general anesthesia）
    \item 術前至少48小時：Aspirin 325 mg QD
    \item 術中需使用 heparin 或 bivalirudin 抗凝
    \item 術後：Aspirin 70-325 mg QD 持續使用
\end{itemize}

\subsection{術者認證要求}

\begin{tcolorbox}[colback=emergency_red!5!white,colframe=emergency_red,title=\textbf{重要：嚴格的術者品質管控}]
\textbf{CAS 術者：}需提交過去12個月的 CAS 案例，並依案例數及結果提交3-25例前瞻性非緊急案例的報告與 angiogram 審查

\textbf{CEA 術者：}需提交前50例連續案例，且 periprocedural stroke/death rate \textbf{<3\%}

$\Rightarrow$ 此嚴格認證確保手術品質，但也意味著研究結果可能無法反映一般臨床實務
\end{tcolorbox}

\subsection{主要終點}

\textbf{Primary outcome：}4年複合終點，包含：
\begin{itemize}
    \item \textbf{Periprocedural period（Day 0-44）：}任何 stroke 或死亡
    \item \textbf{Postprocedural period（Day 44後至4年）：}同側缺血性中風（ipsilateral ischemic stroke）
\end{itemize}

\textbf{Stroke 定義：}依據 WHO 標準，為腦功能局部（或整體）障礙的臨床表現，持續至少24小時，除心血管來源外無其他明顯原因。

% ============================================================
\section{研究結果}
% ============================================================

\subsection{患者基線特徵}

\begin{longtable}{p{5cm}cccc}
\caption{患者基線特徵} \\
\toprule
\multirow{2}{*}{\textbf{特徵}} & \multicolumn{2}{c}{\textbf{Stenting Trial}} & \multicolumn{2}{c}{\textbf{CEA Trial}} \\
\cmidrule(lr){2-3} \cmidrule(lr){4-5}
& \textbf{Medical} & \textbf{Stenting} & \textbf{Medical} & \textbf{CEA} \\
& (n=629) & (n=616) & (n=623) & (n=617) \\
\midrule
\endfirsthead

\toprule
\multirow{2}{*}{\textbf{特徵}} & \multicolumn{2}{c}{\textbf{Stenting Trial}} & \multicolumn{2}{c}{\textbf{CEA Trial}} \\
\cmidrule(lr){2-3} \cmidrule(lr){4-5}
& \textbf{Medical} & \textbf{Stenting} & \textbf{Medical} & \textbf{CEA} \\
\midrule
\endhead

\bottomrule
\endfoot

年齡 (歲) & 69.7±7.7 & 69.3±8.1 & 70.4±7.6 & 70.7±7.8 \\
女性 (\%) & 38.2 & 36.9 & 39.0 & 35.3 \\
白人 (\%) & 90.1 & 92.9 & 88.3 & 90.0 \\
\midrule
\multicolumn{5}{l}{\textbf{危險因子}} \\
高血壓 (\%) & 87.4 & 88.0 & 84.9 & 85.1 \\
糖尿病 (\%) & 37.8 & 40.7 & 38.0 & 34.4 \\
血脂異常 (\%) & 93.3 & 92.0 & 90.0 & 91.5 \\
目前吸菸 (\%) & 21.0 & 18.8 & 21.2 & 21.1 \\
心血管疾病/CABG史 (\%) & 54.5 & 53.7 & 43.3 & 44.3 \\
\midrule
\multicolumn{5}{l}{\textbf{基線數值}} \\
收縮壓 (mmHg) & 138.8±20.2 & 138.2±20.2 & 137.9±19.8 & 139.3±20.2 \\
LDL cholesterol (mg/dL) & 76.7±34.6 & 77.1±36.5 & 81.3±33.9 & 80.3±33.6 \\
BMI (kg/m²) & 28.7±5.6 & 29.3±5.5 & 28.5±5.4 & 28.7±5.3 \\
\midrule
\multicolumn{5}{l}{\textbf{狹窄程度}} \\
$\geq$70\% 狹窄 (\%) & 97.6 & 97.7 & 97.1 & 97.4 \\
PSV $\geq$389 cm/sec (\%) & 33.5 & 31.1 & 32.9 & 37.2 \\
對側 $\geq$50\% 狹窄 (\%) & 34.4 & 37.0 & 37.2 & 35.7 \\
\bottomrule
\end{longtable}

\subsection{危險因子控制情況}

\begin{itemize}
    \item 各組患者的血壓及 LDL cholesterol 控制率在隨機分配後數月內均顯著提升，並維持穩定
    \item 44天時，約 \textbf{60\%} 患者血壓達標（<130 mmHg）
    \item 44天時，約 \textbf{70\%} 患者 LDL cholesterol 達標（<70 mg/dL）
    \item 各治療組間的危險因子控制情況相似
\end{itemize}

\subsection{治療交叉情況}

\begin{longtable}{p{6cm}cc}
\toprule
\textbf{交叉情況} & \textbf{Stenting Trial} & \textbf{CEA Trial} \\
\midrule
\endfirsthead

\bottomrule
\endfoot

內科治療組後來接受 revascularization & 106 (17\%) & 111 (18\%) \\
\quad - 接受 CAS & 78 & 17 \\
\quad - 接受 CEA & 28 & 94 \\
\midrule
介入組未接受手術 & 41 (7\%) & 24 (4\%) \\
\bottomrule
\end{longtable}

\newpage
\subsection{主要終點結果}

\begin{tcolorbox}[colback=yellow!10!white,colframe=orange!75!black,title=\textbf{主要結果摘要}]

\textbf{\large Stenting Trial：CAS 顯著優於單獨內科治療}

\begin{center}
\begin{tabular}{lcc}
\toprule
& \textbf{Medical Therapy} & \textbf{CAS + Medical} \\
\midrule
4年主要終點事件率 & 6.0\% (95\% CI 3.8-8.3) & 2.8\% (95\% CI 1.5-4.3) \\
\bottomrule
\end{tabular}
\end{center}

\begin{itemize}[noitemsep]
    \item \textbf{絕對風險差異：}3.2 percentage points (95\% CI 0.6-5.9)
    \item \textbf{相對風險：}2.13 (95\% CI 1.15-4.39)
    \item \textbf{P value = 0.02}
    \item \textbf{NNT = 31}（治療31人可預防1例主要終點事件）
\end{itemize}

\vspace{0.3cm}
\hrule
\vspace{0.3cm}

\textbf{\large Endarterectomy Trial：CEA 未達統計顯著效益}

\begin{center}
\begin{tabular}{lcc}
\toprule
& \textbf{Medical Therapy} & \textbf{CEA + Medical} \\
\midrule
4年主要終點事件率 & 5.3\% (95\% CI 3.3-7.4) & 3.7\% (95\% CI 2.1-5.5) \\
\bottomrule
\end{tabular}
\end{center}

\begin{itemize}[noitemsep]
    \item \textbf{絕對風險差異：}1.6 percentage points (95\% CI $-$1.1 to 4.3)
    \item \textbf{相對風險：}1.43 (95\% CI 0.78-2.72)
    \item \textbf{P value = 0.24}（未達統計顯著）
\end{itemize}

\end{tcolorbox}

\subsection{終點事件細分分析}

\begin{longtable}{p{5cm}cccc}
\caption{主要終點各組成部分分析} \\
\toprule
\multirow{2}{*}{\textbf{變項}} & \multicolumn{2}{c}{\textbf{Stenting Trial}} & \multicolumn{2}{c}{\textbf{CEA Trial}} \\
\cmidrule(lr){2-3} \cmidrule(lr){4-5}
& \textbf{Medical} & \textbf{Stenting} & \textbf{Medical} & \textbf{CEA} \\
\midrule
\endfirsthead

\toprule
\textbf{變項} & \textbf{Medical} & \textbf{Stenting} & \textbf{Medical} & \textbf{CEA} \\
\midrule
\endhead

\bottomrule
\endfoot

\multicolumn{5}{l}{\textbf{Periprocedural (Day 0-44): Stroke or Death}} \\
事件數/患者數 & 0/629 & 8/616 & 3/623 & 9/617 \\
事件率 (\%) & 0.0 & 1.3 & 0.5 & 1.5 \\
\midrule
\multicolumn{5}{l}{\textbf{Postprocedural (>44 days): Ipsilateral Ischemic Stroke}} \\
追蹤人年 & 1686 & 1714 & 1761 & 1823 \\
事件數/患者數 & 28/600 & 7/582 & 23/600 & 10/596 \\
年事件率 (\%) & 1.7 & 0.4 & 1.3 & 0.5 \\
相對風險 (95\% CI) & \multicolumn{2}{c}{4.07 (1.78-9.31)} & \multicolumn{2}{c}{2.38 (1.13-5.00)} \\
\bottomrule
\end{longtable}

\begin{tcolorbox}[colback=emergency_blue!5!white,colframe=emergency_blue,title=\textbf{重點解讀}]
\begin{itemize}
    \item \textbf{Periprocedural stroke/death：}CAS 組 1.3\%、CEA 組 1.5\%，均高於單獨內科治療組（0-0.5\%）
    \item \textbf{Postprocedural ipsilateral stroke 年發生率：}
    \begin{itemize}
        \item 內科治療組：1.3-1.7\%/年
        \item Revascularization 組：0.4-0.5\%/年
        \item CAS 和 CEA 均顯著降低術後同側中風風險
    \end{itemize}
    \item \textbf{CAS trial 達到顯著差異的關鍵：}術後長期 ipsilateral stroke 的大幅減少（RR 4.07）足以彌補 periprocedural risk
\end{itemize}
\end{tcolorbox}

\subsection{次要終點與安全性}

\begin{longtable}{p{5cm}cccc}
\caption{嚴重不良事件} \\
\toprule
\multirow{2}{*}{\textbf{事件}} & \multicolumn{2}{c}{\textbf{Stenting Trial}} & \multicolumn{2}{c}{\textbf{CEA Trial}} \\
\cmidrule(lr){2-3} \cmidrule(lr){4-5}
& \textbf{Medical} & \textbf{Stenting} & \textbf{Medical} & \textbf{CEA} \\
\midrule
\endfirsthead

\bottomrule
\endfoot

頸動脈血管重建 & 118 (18.8\%) & 29 (4.7\%) & 131 (21.0\%) & 44 (7.1\%) \\
死亡 & 69 (11.0\%) & 48 (7.8\%) & 60 (9.6\%) & 54 (8.8\%) \\
\bottomrule
\end{longtable}

\textbf{內科治療組接受 revascularization 的常見原因：}
\begin{itemize}
    \item 新發頸動脈症狀（new carotid symptoms）
    \item 狹窄進展（progression of stenosis）
    \item 患者偏好（patient preference）
\end{itemize}

\subsection{次族群分析}

各次族群的治療效果大致一致，包括：
\begin{itemize}
    \item 年齡（<70歲 vs $\geq$70歲；<80歲 vs $\geq$80歲）
    \item 性別
    \item Peak systolic velocity（<342 vs $\geq$342 cm/sec）
    \item 高血壓、糖尿病、血脂異常、吸菸狀態
    \item CHA$_2$DS$_2$-VASc score（0-3 vs $\geq$4）
    \item 同側/對側血管的症狀史
\end{itemize}

% ============================================================
\section{討論與臨床意義}
% ============================================================

\subsection{重要發現}

\begin{enumerate}
    \item \textcolor{emergency_red}{\textbf{CAS + 強化內科治療優於單獨內科治療：}}
    \begin{itemize}
        \item 4年主要終點事件率 2.8\% vs 6.0\%（ARD 3.2\%, P=0.02）
        \item NNT = 31
        \item 雖然 periprocedural stroke/death 為 1.3\%，但長期 ipsilateral stroke 的大幅減少（年發生率 0.4\% vs 1.7\%）足以彌補此風險
    \end{itemize}

    \item \textcolor{emergency_red}{\textbf{CEA + 強化內科治療未顯示統計顯著效益：}}
    \begin{itemize}
        \item 4年事件率 3.7\% vs 5.3\%（P=0.24）
        \item 雖有數值上的差異，但未達統計顯著
        \item 可能原因：CEA 的 periprocedural event rate（1.5\%）略高，影響整體效益
    \end{itemize}

    \item \textcolor{emergency_blue}{\textbf{強化內科治療的效果：}}
    \begin{itemize}
        \item 兩個試驗中，單獨內科治療組的年中風率約 1.3-1.7\%
        \item 較過去歷史資料（約 2-3\%/年）明顯降低
        \item 強調現代內科治療的重要性
    \end{itemize}

    \item \textcolor{emergency_blue}{\textbf{兩種 revascularization 的術後 ipsilateral stroke 率相似：}}
    \begin{itemize}
        \item CAS 組：0.4\%/年
        \item CEA 組：0.5\%/年
        \item 兩者均顯著低於單獨內科治療
    \end{itemize}
\end{enumerate}

\subsection{與其他試驗的比較}

\begin{longtable}{p{3cm}p{11cm}}
\toprule
\textbf{試驗} & \textbf{差異與評論} \\
\midrule
\endfirsthead

\bottomrule
\endfoot

\textbf{SPACE-2} & 未顯示 revascularization 效益，但僅達目標收案數的 16\%，統計檢定力不足；且使用 ECST 狹窄標準（較不嚴格） \\
\midrule
\textbf{ECST-2} & 納入 50-69\% 狹窄患者，其中 1/3 有症狀；revascularization 組僅 10/214 接受 CAS \\
\midrule
\textbf{Oxford Vascular Study} & 觀察性研究，年中風率較 CREST-2 低，但依賴患者就醫才能確認事件；CREST-2 採用較敏感的 MRI（82\%）及標準化追蹤問卷 \\
\bottomrule
\end{longtable}

\subsection{臨床實務建議}

\begin{tcolorbox}[colback=emergency_blue!5!white,colframe=emergency_blue,title=\textbf{臨床建議}]
\begin{enumerate}
    \item \textbf{所有 $\geq$70\% 無症狀頸動脈狹窄患者都應接受強化內科治療：}
    \begin{itemize}
        \item 血壓目標 <130 mmHg
        \item LDL cholesterol <70 mg/dL
        \item 戒菸、體重控制、規律運動
    \end{itemize}

    \item \textbf{考慮 CAS 的情況：}
    \begin{itemize}
        \item 高度（$\geq$70\%）無症狀頸動脈狹窄
        \item 具備經認證、經驗豐富的 interventionist
        \item 患者經充分告知後同意
    \end{itemize}

    \item \textbf{CEA 的角色：}
    \begin{itemize}
        \item 雖然本試驗未顯示 CEA 的統計顯著效益，但數值上仍有降低趨勢
        \item 對於不適合 CAS 的患者（如 tortuous anatomy），CEA 仍為合理選項
        \item 需由經驗豐富的 surgeon 執行（periprocedural stroke/death <3\%）
    \end{itemize}

    \item \textbf{選擇單獨內科治療的情況：}
    \begin{itemize}
        \item 患者偏好非侵入性治療
        \item 預期壽命有限
        \item 高手術/介入風險
    \end{itemize}

    \item \textbf{追蹤建議：}
    \begin{itemize}
        \item 定期 carotid duplex ultrasound 監測狹窄進展
        \item 衛教患者 TIA 症狀，若出現應緊急就醫
    \end{itemize}
\end{enumerate}
\end{tcolorbox}

% ============================================================
\section{研究限制}
% ============================================================

\begin{enumerate}
    \item \textbf{非完全雙盲：}雖然 stroke adjudicators 對治療分配不知情，但患者與治療醫師知曉分組

    \item \textbf{試驗期間內科治療持續進步：}
    \begin{itemize}
        \item 血壓目標從 <140 mmHg 調降至 <130 mmHg
        \item PCSK9 inhibitors 的引入
        \item 新型糖尿病/肥胖藥物（如 GLP-1 receptor agonists）的廣泛使用
    \end{itemize}

    \item \textbf{術者選擇偏差：}僅納入經認證的高品質術者，結果可能無法反映一般臨床實務

    \item \textbf{未納入 transcarotid artery revascularization (TCAR)：}此技術在試驗後期才普及

    \item \textbf{Tipping-point analysis 敏感度：}
    \begin{itemize}
        \item Stenting trial：僅需改變 3-4 例事件即可能影響顯著性
        \item 但這代表 >10\% 的相對事件數變化
    \end{itemize}

    \item \textbf{部分術後中風可能與頸動脈狹窄無因果關係：}Revascularization 無法預防所有機轉的中風
\end{enumerate}

% ============================================================
\section{結論}
% ============================================================

\begin{enumerate}
    \item \textcolor{emergency_red}{\textbf{CAS + 強化內科治療}}對於高度無症狀頸動脈狹窄患者，在預防4年內中風/死亡方面，\textbf{顯著優於}單獨強化內科治療（2.8\% vs 6.0\%, P=0.02, NNT=31）

    \item \textcolor{emergency_red}{\textbf{CEA + 強化內科治療}}相較單獨強化內科治療，\textbf{未達統計顯著效益}（3.7\% vs 5.3\%, P=0.24）

    \item \textbf{強化內科治療}是所有患者的治療基礎，可將年中風率控制在 1.3-1.7\%

    \item 此結果挑戰了傳統以 CEA 為無症狀頸動脈狹窄首選治療的觀念

    \item 持續進行的長期追蹤研究將提供更多證據
\end{enumerate}

\vspace{0.5cm}
\textbf{關鍵詞：}Carotid stenosis、Carotid artery stenting、Carotid endarterectomy、Asymptomatic、Intensive medical management、Stroke prevention、CREST-2

% ============================================================
\section{參考文獻}
% ============================================================

\begin{enumerate}
    \item Brott TG, Howard G, Lal BK, et al. Medical Management and Revascularization for Asymptomatic Carotid Stenosis. \href{https://doi.org/10.1056/NEJMoa2508800}{\textit{N Engl J Med}. 2026;394:219-31.}

    \item Brott TG, Hobson RW II, Howard G, et al. Stenting versus endarterectomy for treatment of carotid-artery stenosis. \href{https://doi.org/10.1056/NEJMoa0912321}{\textit{N Engl J Med}. 2010;363:11-23.}

    \item Halliday A, Bulbulia R, Bonati LH, et al. Second asymptomatic carotid surgery trial (ACST-2): a randomised comparison of carotid artery stenting versus carotid endarterectomy. \href{https://doi.org/10.1016/S0140-6736(21)01910-3}{\textit{Lancet}. 2021;398:1065-73.}

    \item Reiff T, Eckstein H-H, Mansmann U, et al. Carotid endarterectomy or stenting or best medical treatment alone for moderate-to-severe asymptomatic carotid artery stenosis: 5-year results of a multicentre, randomised controlled trial. \href{https://doi.org/10.1016/S1474-4422(22)00290-3}{\textit{Lancet Neurol}. 2022;21:877-88.}

    \item White CJ, Brott TG, Gray WA, et al. Carotid artery stenting: JACC state-of-the-art review. \href{https://doi.org/10.1016/j.jacc.2022.05.007}{\textit{J Am Coll Cardiol}. 2022;80:155-70.}
\end{enumerate}

% ============================================================
\section{縮寫對照表}
% ============================================================

\begin{longtable}{p{4cm}p{10cm}}
\toprule
\textbf{縮寫} & \textbf{全名} \\
\midrule
\endfirsthead

\toprule
\textbf{縮寫} & \textbf{全名} \\
\midrule
\endhead

\bottomrule
\endfoot

ARD & Absolute Risk Difference (絕對風險差異) \\
CABG & Coronary Artery Bypass Grafting (冠狀動脈繞道手術) \\
CAS & Carotid Artery Stenting (頸動脈支架置放術) \\
CEA & Carotid Endarterectomy (頸動脈內膜切除術) \\
CI & Confidence Interval (信賴區間) \\
CTA & Computed Tomographic Angiography (電腦斷層血管攝影) \\
ICA & Internal Carotid Artery (內頸動脈) \\
CCA & Common Carotid Artery (總頸動脈) \\
LDL & Low-Density Lipoprotein (低密度脂蛋白) \\
MRA & Magnetic Resonance Angiography (磁振血管攝影) \\
MRI & Magnetic Resonance Imaging (磁振造影) \\
NIHSS & National Institutes of Health Stroke Scale \\
NNT & Number Needed to Treat (需治療人數) \\
PCSK9 & Proprotein Convertase Subtilisin-Kexin Type 9 \\
PSV & Peak Systolic Velocity (收縮期峰值流速) \\
RR & Relative Risk (相對風險) \\
TCAR & Transcarotid Artery Revascularization \\
TIA & Transient Ischemic Attack (暫時性腦缺血發作) \\
WHO & World Health Organization (世界衛生組織) \\
\bottomrule
\end{longtable}

\end{document}
